\subsection{Larger data: scale 4}
\label{sec:scale4}

At scale 4 the dataset has \M{incidencesC} incidences. After loading, HStore's files take \M{hCdisk} and HyperGraphDB's \M{gCdisk}, so both are larger than their 64\,MiB caches, HyperGraphDB's by a much wider margin. HStore's leads hold or grow. Hyperedge insertion becomes \M{sumCingestedges} faster (\M{hCingestedges} against \M{gCingestedges} per second), batched updates \M{sumCwriteupdate} and single-update commits \M{sumClatencycommit}. Co-membership counts are \M{sumCreadcomembership} faster, incidence lookups \M{sumCreadincidence}, and cold lookups after a reopen \M{sumCreadincidencecold}.

Member enumeration is \M{sumCreadmembers} faster: \M{hCreadmembers} hyperedges per second against HyperGraphDB's \M{gCreadmembers}. Two-hop neighbourhoods, which enumerate the members of every incident hyperedge, are \M{sumCreadtwohop} faster. During \wl{read.members}, HStore read nothing from its files and allocated \M{allocChstorereadmembers} per lookup, so every catalog leaf the probes reached was already in its node cache. A cached catalog leaf holds its page bytes and only the records that have been read (\cref{sec:systems}), which keeps the catalog within the 64\,MiB budget. HyperGraphDB serves the same lookups from its atom cache, which holds each link's targets as a Java array and is not limited by the 64\,MiB budget. It allocates \M{allocChgdbreadmembers} per member lookup. A two-hop query allocates \M{allocChstorereadtwohop} in HStore and \M{allocChgdbreadtwohop} in HyperGraphDB.

After the deletes, HStore's files grow to \M{hCdiskchurned}. The deletes rewrite a large part of the incidence index, and the replaced pages remain reachable from the 64 retained generations until the next restart releases that history. The benchmark restarts the store after the deletes, and during the 10\,s \wl{compaction.baseline} window that follows, HStore's background maintenance released and compacted that space. When the measured compaction window began, the explicit compaction took \M{compmsChstore} and found almost nothing left to reclaim, and the final size was \M{hCdiskcompacted}. JE's cleaner ran for \M{compmsChgdb} in the same window and brought HyperGraphDB from \M{gCdiskchurned} to \M{gCdiskcompacted}.

A clean reopen is \M{sumCreopen} faster in HStore, and opening after a crash takes \M{hCrecover} against \M{gCrecover}. Over the whole run, HStore wrote \M{totalwrittenChstore} and HyperGraphDB \M{totalwrittenChgdb}, most of it during cleaning. HStore's single-update commits are faster at every percentile: \M{latChstorelatencycommitptwo} against \M{latChgdblatencycommitptwo} at p99, and \M{latChstorelatencycommitpthree} against \M{latChgdblatencycommitpthree} at p99.9. The median of the per-run maxima is also lower in HStore, \M{latChstorelatencycommitmax} against \M{latChgdblatencycommitmax}, but the slowest single commit across all runs was HStore's: \M{latChstorelatencycommitworst}, against \M{latChgdblatencycommitworst} for HyperGraphDB.

\subsection{Synchronous commits}
\label{sec:sync}

With synchronous commits, every acknowledged commit has been forced to stable storage. On this machine a forced write costs about 4\,ms, and it dominates single-update transactions in both systems. HStore commits \M{hBlatencycommit} of them per second and HyperGraphDB \M{gBlatencycommit} (\M{sumBlatencycommit}). The medians are \M{latBhstorelatencycommitmed} and \M{latBhgdblatencycommitmed}, one forced write each, and HStore's lead comes from the tail: \M{latBhstorelatencycommitpthree} against \M{latBhgdblatencycommitpthree} at p99.9.

Batched work spreads the forced write over 1,000 operations. HStore inserts nodes \M{sumBingestnodes} as fast as HyperGraphDB and hyperedges \M{sumBingestedges} as fast (ranges overlapping), and updates values \M{sumBwriteupdate} as fast. The open-loop writer of \wl{mixed} commits 10 updates every 10\,ms. Both systems held that schedule in every run. Measured from their scheduled start, HStore's commits took \M{latBhstoremixedwritemed} at the median and \M{latBhstoremixedwriteptwo} at p99, and HyperGraphDB's \M{latBhgdbmixedwritemed} and \M{latBhgdbmixedwriteptwo}.

Under synchronous commits, JE reports \M{bpoBhgdblatencycommit} written per single-update commit, against \M{bpoAhgdblatencycommit} with asynchronous commits. The extra bytes come with the log flush that the adapter adds (\cref{sec:durability}). We have not established what JE writes in that flush, so the bytes-written figures for synchronous single commits in \cref{app:full} should not be read as a comparison.

Reads do not commit, and HStore leads every read workload in this configuration as well at the median, with overlapping ranges for member enumeration and giant-edge scans (\cref{tab:summary}). Recovery after SIGKILL takes \M{hBrecover} in HStore and \M{gBrecover} in HyperGraphDB at the median, and the ranges do not overlap. As in every other configuration, no run of either system lost an acknowledged commit.

